% Generated draft: requires mathematical and source review.
\documentclass[11pt]{article}
\usepackage[margin=25mm]{geometry}
\usepackage{fontspec}
\setmainfont{Latin Modern Roman}
\usepackage{amsmath,amssymb,mathtools,unicode-math}
\setmathfont{Latin Modern Math}
\usepackage{enumitem,needspace,xurl,hyperref}
\setlength{\parindent}{0pt}
\setlength{\parskip}{5pt}
\setlength{\emergencystretch}{3em}
\newcommand{\N}{\mathbb{N}}
\newcommand{\Z}{\mathbb{Z}}
\newcommand{\Q}{\mathbb{Q}}
\newcommand{\R}{\mathbb{R}}
\newcommand{\C}{\mathbb{C}}
\newcommand{\F}{\mathbb{F}}
\newcommand{\A}{\mathbb{A}}
\newcommand{\PP}{\mathbb{P}}
\newcommand{\E}{\mathbb{E}}
\newcommand{\eps}{\varepsilon}
\newcommand{\dd}{\,\mathrm{d}}
\begin{document}
\section*{UnsolvedMath ID: 30002422}
\textbf{Source problem code:} catalogue entry
\subsection{Definitions and mathematical statement}
Let \(Q=[0,1]^2\), let \(|E|\) denote planar Lebesgue measure, and write \(E\mathbin\triangle F=(E\setminus F)\cup(F\setminus E)\). For \(N\in\N\), consider a measurable partition \(\mathcal P=(\Omega_1,\ldots,\Omega_N)\) of \(Q\), with cells disjoint up to sets of measure zero. Require \(|\Omega_i|>0\) for every \(i\); this restriction is necessary for the displayed ratios, although it is not explicit in the supplied statement. Define
\[
\mathcal A(\Omega_i)=\inf_{\substack{B\subset\R^2\text{ a disk}\\|B|=|\Omega_i|}}\frac{|\Omega_i\mathbin\triangle B|}{|\Omega_i|},\qquad
\mathcal D(\Omega_i)=\frac{|\Omega_i|-\min_{1\le j\le N}|\Omega_j|}{|\Omega_i|},
\]
and
\[
F(\mathcal P)=\sum_{i=1}^{N}|\Omega_i|\mathcal A(\Omega_i)+\sum_{i=1}^{N}|\Omega_i|\mathcal D(\Omega_i).
\]
Here the comparison disks are interpreted as arbitrary planar disks; the supplied record imposes no containment requirement. Let \(\mathfrak P_N\) be this class of positive-area partitions and put
\[
m_N=\inf_{\mathcal P\in\mathfrak P_N}F(\mathcal P).
\]
Using an infimum does not assert existence of a minimizing partition.

\textbf{Source question.} Is the extremal partition for the stated geometric uncertainty principle a decomposition of \(Q\) into hexagons, with boundary modifications? The supplied expanded formulation asks whether the minimum of \(F\) is asymptotically attained by such partitions.

\textbf{Unresolved qualifications.} The record does not specify the asymptotic parameter, the geometry and sizes of the hexagons, or the admissible boundary modifications. If the intended regime is \(N\to\infty\) and asymptotic attainment means vanishing absolute error, a possible precise target is the existence of partitions \(\mathcal H_N\in\mathfrak P_N\) of the prescribed hexagonal type such that
\[
\lim_{N\to\infty}\bigl(F(\mathcal H_N)-m_N\bigr)=0.
\]
This is a conditional interpretation requiring review, not a uniquely determined formulation supplied by the record. A definition of the hexagonal type and its boundary modifications is still needed. The original wording about an extremal partition also leaves unresolved whether exact finite-partition optimality was intended.
\subsection{Short English statement}
The question asks whether hexagonal decompositions of the unit square, suitably adjusted at its boundary, minimize asymptotically a weighted sum of distance from disks and inequality of cell areas. Positive cell areas are needed to define the functional. The limiting regime and the admissible hexagonal constructions remain unspecified.
\subsection{Sources}
\begin{itemize}
\item Local UnsolvedMath problem registry, permanent ID \(30002422\), title: \emph{Hexagonal Extremizers for a Geometric Uncertainty Principle}; exact source problem code: \texttt{catalogue entry}. Collection: Oberwolfach Reports Open Problems. No individual proposer is supplied; the proposed year is recorded as 2013.
\item Supplied report citation: \emph{Uniform Distribution Theory and Applications} (2014), Oberwolfach Reports, pp.\ 2837--2917, DOI \texttt{10.4171/owr/2013/49}. Supplied source URL and literature-source link, labeled ``original OWR report'': \url{https://doi.org/10.4171/owr/2013/49}. The displayed citation year is retained as supplied; its relation to the 2013 proposal year and DOI has not been independently checked.
\item Supplied literature assessment, checked on \textbf{2026-08-21}: status recorded as open; the hexagonal decomposition is described as a conjectural extremizer for the quantitative partition inequality, and the curation reports that no matching optimality theorem was found. This is supplied dated curation, not an independently verified literature assessment. The background explicitly requests independent verification of time-sensitive status and cited claims.
\item The registry labels the statement \texttt{corrected\_verified} and reports that definitions of Fraenkel asymmetry, relative size deviation, and the weighted partition functional were restored from the official report. This verification claim is supplied metadata and has not been independently checked here.
\end{itemize}
\end{document}
